\documentclass[11pt]{article}
\usepackage[a4paper,margin=1in]{geometry}
\usepackage{amsmath,amssymb,mathtools,bm}
\usepackage{enumitem,booktabs,array}
\usepackage{tcolorbox}
\usepackage{hyperref}
\usepackage[protrusion=true,expansion=false]{microtype}
\hypersetup{colorlinks=true,linkcolor=blue,urlcolor=blue,citecolor=blue}
\setlist[itemize]{topsep=2pt,itemsep=2pt}
\setlist[enumerate]{topsep=2pt,itemsep=2pt}
\newtcolorbox{keybox}{colback=blue!3!white,colframe=blue!40!black,boxrule=0.4pt,arc=1mm}
\newtcolorbox{alertbox}{colback=red!3!white,colframe=red!40!black,boxrule=0.4pt,arc=1mm}
\newtcolorbox{answerbox}{colback=green!3!white,colframe=green!40!black,boxrule=0.4pt,arc=1mm}
\newtcolorbox{drillbox}{colback=yellow!5!white,colframe=orange!60!black,boxrule=0.4pt,arc=1mm}
\newcommand{\EV}{\mathbb{E}}
\newcommand{\Var}{\mathrm{Var}}
\newcommand{\Cov}{\mathrm{Cov}}
\newcommand{\blank}[1]{\underline{\hspace{#1}}}

\title{W14-03: Keys \& Mulder (2025, WP) \\
\large ``Neglected No More: Housing Markets, Mortgage Lending, and Sea Level Rise'' --- Active Recall Workbook}
\author{Banking \& Risk Management --- Exam Prep}
\date{\today}
\begin{document}\maketitle

\begin{keybox}
\textbf{Paper in one sentence.} Using 2011--2023 US housing transactions and mortgage data, KM show that the SLR discount in coastal home prices has \emph{grown} to $\sim$5--10\%, mortgage originations in SLR-exposed areas have slowed, and local banks charge higher rates --- so housing and mortgage markets are now adjusting to SLR risk, reversing Murfin-Spiegel's (2020) null.

\textbf{Placement.} Dynamic update to the BGL vs.\ MS debate. Documents that climate-risk pricing evolves as information and attention improve; a crucial paper for the ``is climate risk priced?'' literature.
\end{keybox}

\section*{Round 1: Complete Walk-Through}

\subsection*{1.1 Data and design}
Housing transactions 2011--2023, extended past the MS sample. Mortgage data from HMDA (originations, loan terms). NOAA SLR inundation classifiers at property level. Hedonic within-ZIP-quarter specifications plus difference-in-differences using information shocks (major climate reports, FEMA updates).

\subsection*{1.2 Headline pricing results}
\begin{itemize}
\item Full-sample SLR discount: $\sim$5--10\% (larger than BGL 2019 on later data).
\item Time trend: discount grows 2011--2023, not flat.
\item Still heterogeneous: high-belief coastal counties discount more (consistent with BBGL 2022).
\end{itemize}

\subsection*{1.3 Mortgage-market results}
\begin{itemize}
\item Origination volume in SLR-exposed areas grows more slowly than control areas.
\item Loan rates modestly higher.
\item Non-local lenders exit; local lenders remain but with more conservative underwriting.
\end{itemize}

\subsection*{1.4 Mechanism}
\begin{itemize}
\item Improved information (climate reports, news coverage, FEMA flood-map updates) shifts buyer and lender beliefs.
\item Mortgage market is slower to adjust than housing but now exhibits clear pricing.
\item Insurance markets are a parallel pressure point; NFIP premium increases feed through.
\end{itemize}

\subsection*{1.5 Implications}
\begin{itemize}
\item Reconciles BGL 2019, MS 2020, BBGL 2022, and newer data: discount is real and growing.
\item Climate-risk pricing is dynamic; early-period null does not imply permanent null.
\item Financial-stability implications: coastal-heavy banks face rising mortgage and property-collateral risk.
\end{itemize}

\subsection*{1.6 Caveats}
\begin{alertbox}
\begin{itemize}
\item Estimation uses NOAA SLR mapping; alternative physical projections may give different exposures.
\item Insurance-market effects are partly co-determined with housing, so mortgage effects may reflect insurance pass-through.
\item Sample includes COVID-era housing boom, which could amplify measured changes.
\end{itemize}
\end{alertbox}

\section*{Round 2: Level 1 Fill-in}
\begin{enumerate}
\item Sample period: \blank{1cm}--\blank{1cm}.
\item SLR discount: $\sim$\blank{1cm}--\blank{1cm}\%.
\item Trend: \blank{1cm} (flat/growing).
\item Mortgage origination in exposed areas: \blank{1cm} (faster/slower) growth.
\item Reconciles with: Murfin--\blank{1cm} (2020).
\end{enumerate}

\section*{Round 3: Level 2 Prompts}
\begin{enumerate}
\item Write the hedonic spec and discuss why 2011--2023 data yield different estimates than 2007--2016.
\item Explain the mortgage-market evidence as evidence on lender beliefs.
\item Discuss financial-stability implications of rising coastal-mortgage risk.
\item How does KM reconcile with BGL, MS, and BBGL?
\end{enumerate}

\section*{Round 4: Minimal Scaffolding}
From a blank page: (i) data, (ii) pricing results, (iii) mortgage results, (iv) mechanism, (v) implications.

\section*{Round 5: Blank Page}
Essay: SLR risk is being priced --- what the housing and mortgage markets now show.

\vspace{6pt}
\begin{drillbox}
\textbf{Speed Drill.} (1) Sample years. (2) Discount magnitude. (3) Mortgage origination direction. (4) Belief-information mechanism. (5) Reconciles which prior paper?
\end{drillbox}

\section*{Quick Reference \& Answer Key}
\begin{answerbox}
\begin{itemize}
\item \textbf{Sample.} 2011--2023 housing + HMDA.
\item \textbf{Discount.} $\sim$5--10\%, growing.
\item \textbf{Mortgages.} Slower origination in exposed areas.
\item \textbf{Reconciles.} Murfin--Spiegel 2020's null.
\end{itemize}
\end{answerbox}

\begin{keybox}
\textbf{30-second exam pitch.} Keys \& Mulder (2025) update the climate-and-housing debate with 2011--2023 transaction and HMDA mortgage data. They find a sea-level-rise discount of $\sim$5--10\% in coastal home prices that has \emph{grown} over time, mortgage origination volume growing more slowly in SLR-exposed areas, higher loan rates, and a shift from non-local to local lenders who underwrite more conservatively. The discount is larger in high-climate-belief counties (Bernstein-Billings-Gustafson-Lewis 2022). The paper reconciles the BGL-MS debate: MS's 2020 null reflected an early sample when information was limited and subsidence dominated; as information has improved and beliefs have shifted, SLR risk is now visibly priced in both housing and mortgage markets. Financial-stability implication: coastal-heavy banks face rising mortgage and collateral risk as climate-risk pricing matures.
\end{keybox}

\end{document}
